%=============================================================================
\section{Introduction}
\label{sec:intro}
%=============================================================================

Consider two bottled water brands in 2030. Both comply with the EU mandate
requiring 30\% recycled plastic content, and both face the same higher material
cost. One sells primarily to environmentally engaged urban consumers; the other
to price-sensitive rural households. The regulation and the material are
identical, yet consumer reactions may differ sharply. Compliance is
enforceable. What is not known is which segments will accept, resist, or
switch, and that knowledge is needed before market outcomes make it observable.

This is the \emph{policy-before-data} problem: regulatory timelines precede
behavioural measurement, so firms and regulators must decide on pricing,
labelling, communication, and sequencing while the relevant consumer responses
remain unmeasured in the target population. It recurs across sustainability
transitions, where mandates are designed to drive structural change before
market actors have adapted \citep{mao2020}. Regulation (EU)~2025/40 is a current
instance. Article~7 sets minimum recycled-content shares of 30\% for
contact-sensitive PET packaging and for single-use plastic beverage bottles
from 2030, rising to 50\% and 65\% respectively from 2040 \citep{eu2025}. The
Regulation entered into force in February~2025 and applies from August~2026.

Three features of this case define the governance problem. First,
\emph{compliance does not secure acceptance}. A package can meet every legal
requirement and still be distrusted \citep{polyportis2022}. If resistant
segments account for most consumption in a product category, a uniform
implementation may erode the legitimacy that later, more ambitious mandates
will need \citep{molenbuur2025}. Second, \emph{response is heterogeneous}.
Reactions to recycled content vary with age, education, environmental concern,
and recycling behaviour \citep{ruokamo2022,zwicker2020}, so an aggregate
acceptance rate hides where the risk lies. Third, \emph{the evidence gap is
built into the timeline}. The Spanish dataset used here, fielded in 2023, does
not measure acceptance of post-consumer recycled (PCR) plastic in food-contact
applications, willingness to pay a premium for it, or its perceived safety and
quality. A dedicated national survey would take months and return a single
static snapshot.

Existing instruments cover only part of this problem. Delphi panels,
cross-impact analysis, and scenario planning structure expert judgement when
quantitative evidence is scarce, but they yield a small number of narrative
futures and are difficult to update \citep{bas2015}. Statistical forecasting
and Bayesian demand updating \citep{lee2012} provide uncertainty estimates but
assume that the target variable has been observed in the population of
interest. Synthetic-data generators based on Bayesian networks reproduce survey
structure with explicit uncertainty \citep{pearl1988,kjaerulff2008,martins2024},
yet they have been developed for disclosure control on a single dataset, for a
small number of binary variables, and for statistically trained users.
Section~\ref{sec:theory} shows that this gap is not incidental: anticipatory
governance and risk-mitigation theory imply specific requirements that none of
these instruments meets jointly.

\paragraph{Aim and contributions.}
This paper examines how an AI-enabled probabilistic decision-support system can
support the anticipatory governance of a mandated circular-packaging transition
when target-specific behavioural evidence is unavailable, and under which
conditions its outputs may be used. It makes three contributions.

\emph{Theoretical.} We frame policy-before-data as a demand-side Collingridge
dilemma and derive from anticipatory governance and risk-mitigation theory four
design requirements for an instrument operating in the pre-observational window
(Section~\ref{sec:theory}).

\emph{Technological.} We present Bayesian Virtual Surveys (BVS), a system
combining language-model-assisted evidence extraction, a Bayesian inference
engine that generates national microdata as in \citet{martins2024} and attaches
unmeasured outcomes through a documented transfer model updated by conjugate
fusion as national data arrive, and a
natural-language interface for non-specialist users
(Sections~\ref{sec:system}--\ref{sec:framework}). The language model extracts
evidence and translates queries; it does not generate acceptance
probabilities.

\emph{Empirical and governance-related.} Applied to Spain under
Regulation~2025/40, the system produces segment-level scenarios for outcomes
not yet measured nationally. We separate what has been validated (reproduction
of observed variables) from what remains a prior-informed projection (the
target outcome), report where external evidence fails to transfer, and
translate the scenarios into a risk-mitigation procedure based on pilots,
engagement, measurement, and updating
(Sections~\ref{sec:results}--\ref{sec:implications}).

